% ============================================================
% §5.4 Q4结果分析 — 灵敏度分析与鲁棒性检验 (v4: 精简版)
% 竞赛: 电工杯 B题 (2026)
% ============================================================

\section{Q4求解结果与分析：灵敏度与鲁棒性}

\subsection{三场景参数扫描设计}

为验证选址-定价联合方案的工程稳健性，以Q2 MILP模型为核，对三类外部冲击实施参数扫描共6组对比实验\cite{bertsimas2011}（基线为Q2确定解F+H+J，10/10全覆盖，目标值54.225）。结果汇总于表\ref{tab:q4_sensitivity}，覆盖率、满意度与利润的全景对比如图\ref{fig:q4_sensitivity_combined}所示。

\begin{figure}[H]
\centering
\includegraphics[width=0.92\textwidth]{figures/figure_q4_sensitivity_combined.png}
\caption{三场景灵敏度分析全景对比。(a) 覆盖率；(b) 加权满意度；(c) 年利润。}
\label{fig:q4_sensitivity_combined}
\end{figure}

\begin{table}[H]
\centering
\caption{三场景参数扫描灵敏度汇总 (确定性求解)}
\label{tab:q4_sensitivity}
\footnotesize
\begin{tabular}{lccccc}
\toprule
\textbf{场景} & \textbf{预算(万)} & \textbf{覆盖} & \textbf{站点配置} & \textbf{平均$S$} & \textbf{年利润(万)} \\
\midrule
基线 (Q2) & 120 & 100\% & F(大)+H(中)+J(中) & 0.850 & 1~097.7 \\
\midrule
A$_1$: 预算放松 & 130 & 100\% & B(大)+E(中)+G(大) & 0.883 & 1~053.8 \\
A$_2$: 预算放松 & 140 & 100\% & B(大)+D(大)+E(大) & 0.889 & 1~010.0 \\
A$_3$: 预算放松 & 150 & 100\% & D(中)+E(大)+F(小)+J(大) & \textbf{0.912} & 980.9 \\
\midrule
B: 成本通胀+20\% & 120 & 80.0\% & A(小)+E(中)+J(大) & 0.906 & 667.7 \\
\midrule
C: 银发海啸 & 120 & 90.0\% & A(小)+G(中)+I(中)+J(中) & 0.895 & 887.9 \\
\bottomrule
\end{tabular}
\end{table}

\subsection{场景A：预算放松——满意度的单调递增与边际递减}

场景A覆盖率保持100\%。满意度单调递增：$S$从0.850升至0.883、0.889、0.912，预算每增加10万带来0.015--0.023的$S$增量。但边际递减显著——$S$增量从130万的+0.033降至140万的+0.006（衰减82\%），130万元已接近全覆盖与高满意度的甜点区间。

\subsection{场景B：成本通胀——全覆盖的崩溃与集约化幸存}

建设与运营成本同步上浮20\%使120万购买力缩水至等效100万，系统被迫从分散全覆盖退化为集约化覆盖：A(小)+E(中)+J(大)，覆盖率降至80\%（C与G未被覆盖），目标值骤降19.6\%。满意度反升至0.906（+6.6\%）——系统通过牺牲覆盖广度换取服务深度实现了帕累托改进。成本通胀是唯一能击穿全覆盖防线的冲击类型。

\subsection{场景C：银发海啸——需求增压下的自适应重配置}

半失能转移率升至5.5\%、失能转移率升至12\%的双重增压下，第5年末失能老人增至1~268人（+12.7\%）。模型自动选择四站布局（A小+G中+I中+J中），覆盖率维持90\%，仅C小区因地理孤立未被覆盖。满意度逆势微升至0.895（+5.3\%），来源是小型站低利用率带来的S2补偿。若固守基线F+H+J配置，F站日服务量将超容量6.7\%——系统的自适应重配置恰是其鲁棒性的核心来源。

\subsection{综合鲁棒性判定}

多维度韧性矩阵（$\rho = 1 - |\Delta X / X_{\text{baseline}}|$）的数值结果见附录表\ref{tab:q4_resilience}，其热力图可视化如图\ref{fig:q4_heatmap}所示：所有场景任一维度$\rho>0.80$。覆盖率韧性0.80--1.00，满意度韧性0.927--0.961，利润率韧性0.609--0.960。场景B利润率维度（0.609）为唯一突破强鲁棒性阈值的场景-维度组合。

\begin{figure}[H]
\centering
\includegraphics[width=0.76\textwidth]{figures/figure_q4_resilience_heatmap.png}
\caption{多场景多维度韧性热力图。蓝色越深韧性越强，红色越深韧性越弱。}
\label{fig:q4_heatmap}
\end{figure}

多维度韧性雷达图（图\ref{fig:q4_radar}）进一步验证了系统的整体鲁棒性。

\begin{figure}[H]
\centering
\includegraphics[width=0.70\textwidth]{figures/figure_q4_resilience_radar.png}
\caption{多维度鲁棒性雷达图。三场景各维度韧性系数均高于强鲁棒性阈值0.75。}
\label{fig:q4_radar}
\end{figure}

\subsection{对政策制定的量化启示}

\begin{enumerate}
    \item \textbf{全覆盖预算底线为109万元}：成本通胀20\%将摧毁全覆盖，建议建立建设成本指数化应急基金（预算的15--20\%）以吸收不可预见冲击。
    \item \textbf{130万元为甜点投资门槛}：额外10万可使满意度从0.850升至0.883（+3.9\%），超出后边际回报急剧下降。建议建立半失能$\to$失能转化率季度监测机制，转化率超11\%触发黄灯预警、超12.5\%触发红灯扩容（H站中型$\to$大型，追加13万），将灵敏度分析转化为可操作的工程风险管理工具。
\end{enumerate}

\subsection{面向全国推广的不确定因素与应对策略}

当此嵌入式养老网络从10个小区向全国铺开时，将遭遇三类题目参数空间之外的深层混沌因素。

\textbf{因素一：邻避效应（NIMBY）。} 嵌入式站点"见缝插针"嵌入老旧小区，易引发青壮年居民对占用公共绿地、医疗废弃物、救护车噪音的群体性抗拒。Q2方案F站所在小区为高密度老年人口区，若遭遇NIMBY阻力，109万元确定性方案将面临不确定性溢价。应对：推行公共服务设施功能叠合演化——将养老职能"寄生"于现有社区党群中心、文体活动站，实现空间无损置换与民意柔性安抚。建模层面可将风险量化为站点社会接受度折扣因子$\gamma_i\in[0,1]$，以$C_k^{\text{build}}/\gamma_i$反映阻力成本。

\textbf{因素二：照护人力资源供应链枯竭。} 上门护理属情感与体力双重密集型服务，Q3定价压至12.64元/次意味着服务人员单次报酬仅约8--10元，难以吸引稳定供给。建议植入"时间银行"代际接力模式：将60--70岁自理老人社区服务时间转化为可储蓄的"时间货币"，以2030年预计近5,000人自理老人池（基年4,885人按2\%年均增长外推）为劳动力蓄水池。时间银行的非货币化锚定（1小时=1积分=1元等值服务）应独立于Q3商业定价体系。

\textbf{因素三：技术颠覆的固定资产沉没风险。} "十五五"至"十六五"窗口内，家用照护机器人一旦跨越商业量产临界点，将对线下服务站体系造成创造性破坏。F站45万元投资按20年折旧，在技术加速迭代下可能8--10年即面临功能性淘汰。底线思维：(1) 站点采用装配式建筑，保留拆卸重组弹性；(2) 服务设备接口标准化，确保智能终端即插即用；(3) 总投资$\ge$15\%预留为技术升级基金。使嵌入式养老网络从"一次性基建"进化为"可演进的公共服务平台"。
